\documentclass[11pt]{article}

\usepackage[T1]{fontenc}
\usepackage[utf8]{inputenc}
\usepackage{libertinus}
\usepackage[letterpaper,margin=1in,headheight=15pt]{geometry}
\usepackage{microtype}
\usepackage{amsmath,amssymb,amsthm,mathtools,bm,mathrsfs}
\usepackage{booktabs,longtable,tabularx,array,multirow}
\usepackage{enumitem}
\usepackage{xcolor}
\usepackage{graphicx}
\usepackage{fancyhdr}
\usepackage{titlesec}
\usepackage{csquotes}
\usepackage{xurl}
\usepackage{hyperref}
\usepackage[nameinlink,noabbrev]{cleveref}
\usepackage{listings}
\usepackage{etoolbox}

\definecolor{linkblue}{RGB}{22,72,118}
\definecolor{darkgreen}{RGB}{25,95,60}
\definecolor{darkred}{RGB}{135,35,35}
\definecolor{shade}{RGB}{247,248,250}
\hypersetup{
  colorlinks=true,
  linkcolor=linkblue,
  citecolor=darkgreen,
  urlcolor=linkblue,
  pdftitle={Logarithmic q-Richardson Acceleration and Lambert Phase Locking in Fabius-Rvachev Analysis},
  pdfauthor={OpenAI research synthesis},
  pdfsubject={New frontier results for Fabius, Rvachev up, Thue-Morse, q-binomial coefficients, sinc products, and Lambert W}
}

\pagestyle{fancy}
\fancyhf{}
\fancyhead[L]{\small Logarithmic $q$-Richardson and Lambert phase locking}
\fancyhead[R]{\small Frontier report}
\fancyfoot[C]{\thepage}

\titleformat{\section}{\Large\bfseries}{\thesection}{0.75em}{}
\titleformat{\subsection}{\large\bfseries}{\thesubsection}{0.65em}{}
\titleformat{\subsubsection}{\normalsize\bfseries}{\thesubsubsection}{0.55em}{}

\newtheorem{theorem}{Theorem}[section]
\newtheorem{proposition}[theorem]{Proposition}
\newtheorem{lemma}[theorem]{Lemma}
\newtheorem{corollary}[theorem]{Corollary}
\newtheorem{conjecture}[theorem]{Conjecture}
\theoremstyle{definition}
\newtheorem{definition}[theorem]{Definition}
\newtheorem{algorithm}[theorem]{Algorithm}
\newtheorem{remark}[theorem]{Remark}
\theoremstyle{remark}
\newtheorem{observation}[theorem]{Observation}

\newcommand{\sinc}{\operatorname{sinc}}
\newcommand{\Up}{\operatorname{up}}
\newcommand{\sgn}{\operatorname{sgn}}
\newcommand{\Li}{\operatorname{Li}}
\newcommand{\dd}{\,\mathrm{d}}
\newcommand{\e}{\mathrm{e}}
\newcommand{\ii}{\mathrm{i}}
\newcommand{\R}{\mathbb{R}}
\newcommand{\C}{\mathbb{C}}
\newcommand{\N}{\mathbb{N}}
\newcommand{\Z}{\mathbb{Z}}
\newcommand{\Q}{\mathbb{Q}}
\newcommand{\abs}[1]{\left\lvert #1\right\rvert}
\newcommand{\norm}[1]{\left\lVert #1\right\rVert}
\newcommand{\floor}[1]{\left\lfloor #1\right\rfloor}
\newcommand{\ceil}[1]{\left\lceil #1\right\rceil}
\newcommand{\brak}[2]{\genfrac{[}{]}{0pt}{}{#1}{#2}}
\newcommand{\poch}[2]{(#1;#2)}
\newcommand{\qbinom}[3]{\genfrac{[}{]}{0pt}{}{#1}{#2}_{#3}}
\newcommand{\lambert}{W_{-1}}
\newcommand{\repo}{\textsc{ProveIt}}
\newcommand{\statusproved}{\textcolor{darkgreen}{\textbf{proved here}}}
\newcommand{\statusconditional}{\textcolor{linkblue}{\textbf{proved conditionally on the repository's all-orders expansion}}}
\newcommand{\statusopen}{\textcolor{darkred}{\textbf{open}}}

\setlist[itemize]{leftmargin=1.5em,itemsep=0.25em,topsep=0.35em}
\setlist[enumerate]{leftmargin=1.7em,itemsep=0.25em,topsep=0.35em}
\setlength{\parindent}{1.2em}
\setlength{\parskip}{0.2em}
\setcounter{tocdepth}{2}
\allowdisplaybreaks

\begin{document}
\hypersetup{pageanchor=false}
\pagenumbering{gobble}

\begin{titlepage}
\centering
\vspace*{1.1cm}
{\Huge\bfseries Logarithmic $q$-Richardson Acceleration\\[0.25em]
and Lambert Phase Locking\\[0.4em]
in Fabius--Rvachev Analysis\par}
\vspace{0.8cm}
{\Large Repository-wide audit, exact sinc-tail extrapolation,\\
Bell-polynomial error calculus, and extraction of the periodic endpoint factor\par}
\vspace{1.2cm}
{\large OpenAI research synthesis\par}
\vspace{0.35cm}
{\large 27 August 2026\par}
\vfill
\begin{minipage}{0.88\textwidth}
\small
\textbf{Research status.} The sinc-product results in Parts~I--III are proved from first principles in this report. The Lambert-lattice extrapolation theorem is an algebraic consequence of the all-orders lower-Lambert expansion developed in the repository; its use is therefore explicitly conditional on that expansion with the stated uniformity. The numerical tables were independently recomputed at 100-digit precision. ``New'' means absent as an exact result from the audited repository snapshot; no claim of worldwide priority is made.
\end{minipage}
\vfill
\end{titlepage}
\clearpage
\hypersetup{pageanchor=true}
\pagenumbering{roman}

\begin{abstract}
The Fabius function, Rvachev's compactly supported $C^\infty$ up-function, the signed Thue--Morse sequence, and dyadic $q$-calculus form a tightly connected system. The existing \repo{} corpus already contains extensive exact dyadic arithmetic, Gaussian-binomial interpolation, moment and cumulant formulas involving Bernoulli and Bell polynomials, Fourier products, Lambert-$W$ endpoint asymptotics with a nonconstant periodic factor, and several forms of Richardson extrapolation. After a repository-wide audit of the living documents, frozen historical versions, research drafts, and bundled source papers, this report develops two extrapolation mechanisms that are not stated there.

First, for the generalized lacunary sinc product
\[
P_b(t)=\prod_{k\ge 1}\sinc(t/b^k),\qquad b>1,
\]
we obtain an exact Bernoulli--zeta expansion of the omitted logarithmic tail. Applying Lagrange interpolation to the logarithms of the truncations at the geometric nodes $q^j$, $q=b^{-2}$, gives a multiplicative Richardson transform with closed Gaussian-$q$-binomial weights. Its full error series is exact, has a fixed sign $(-1)^r$ at level $r$, yields consecutive certified upper and lower bounds, and admits an explicit nonasymptotic remainder estimate. The $\ell^1$ norm of the weights is exactly $(-q;q)_r/(q;q)_r$, hence remains bounded as the order tends to infinity; in the dyadic case it stays below $1.970$. The method therefore combines high order with unusually mild conditioning. Bell polynomials convert the logarithmic error into relative-error and derivative formulas, while a finite Thue--Morse exponential-sum identity transfers the accelerated product to purely finite signed digital sums.

Second, the exact inverse coordinate
\[
\lambda(x)=-\frac{1}{\log 2}W_{-1}(-x\log 2),\qquad x=\lambda 2^{-\lambda},
\]
is used to build phase-locked nodes $x_j=(\lambda+j)2^{-(\lambda+j)}$. Periodic coefficient functions are identical at all these nodes. Ordinary Lagrange extrapolation in $h_j=(\lambda+j)^{-1}$ then cancels the first $r$ inverse-$\lambda$ corrections and extracts the periodic endpoint term $\Psi(\lambda)$ directly from values of $F$. The exact weights, all transformed moments, and their complete-homogeneous/Bell-polynomial representation are derived. Bernoulli polynomials describe the arithmetic-node power sums needed in conditioning and preconditioning. Together these results provide a new bridge between $q$-binomial interpolation, sinc products, Thue--Morse blocks, Lambert-$W$ phase geometry, Bernoulli coefficients, Bell-polynomial reconstruction, and Richardson extrapolation.
\end{abstract}

\tableofcontents
\clearpage
\pagenumbering{arabic}

\section{Scope, corpus audit, and nonduplication boundary}

\subsection{The audited corpus}

The target tree contains living expositions, a deliberately quarantined frontier volume, active drafts, frozen historical versions, and source-paper transcriptions. The audit used the current tree and the archive ledger as of 27 August 2026. Counting each distinct \texttt{.tex} path, the inventory contains 57 sources: 25 in the living tree and 32 frozen historical sources. The living set consists of four top-level expositions, the canonical frontier volume, one Thue--Morse atlas, twelve Fourier-decay dossiers, and seven source-paper transcriptions.

\begin{table}[htbp]
\centering
\caption{Repository-wide source inventory used for the overlap audit.}
\label{tab:inventory}
\begin{tabularx}{\textwidth}{@{}l r X@{}}
\toprule
Family & TeX files & Principal content relevant to this report \\
\midrule
Living primary expositions & 4 & Fabius/up synthesis, glossary, non-elementarity, Lean walkthrough \\
Canonical frontier volume & 1 & Dyadic web, repeated integration, endpoint asymptotics, Thue--Morse convergence, saddle analysis \\
Living formula atlas & 1 & Thue--Morse identities, Prouhet cancellation, rarefaction, Mahler structure, Bell/Bernoulli formulas \\
Living Fourier-decay dossiers & 12 & Infinite sinc products, zeros, dyadic renormalization, envelopes, Lambert-$W$ scales, competing asymptotic formulations \\
Bundled source-paper TeX & 7 & Lacunary products, up-function background, arithmetic of $F$, approximation literature \\
Frozen archive & 32 & Earlier notes, parallel expositions, q-calculus articles, supplements, formula atlases, topic dossiers \\
\midrule
Total & 57 & Every path is listed in Appendix~\ref{app:inventory}. \\
\bottomrule
\end{tabularx}
\end{table}

The archive matters for originality. Several ideas no longer visible in the living primary paper survive in an older parallel draft or a superseded dossier. In particular, the audit treated the following as prior art rather than rediscoveries:
\begin{itemize}
\item Gaussian-$q$-binomial coefficient extraction at geometric nodes;
\item exact dyadic Richardson transforms for sequences modeled by powers of $2^{-p}$;
\item Bernoulli-number cumulants and Bell-polynomial moment reconstruction;
\item Prouhet cancellation for signed Thue--Morse blocks;
\item fixed-numerator Lambert phase locking along dyadic endpoint rays;
\item the periodic factor in the corrected small-argument asymptotic;
\item exact finite Thue--Morse products and Bromwich/coefficient representations;
\item global Fourier-decay envelopes and the zero geometry of the up transform.
\end{itemize}

\subsection{What is new relative to the repository}

The exact constructions below were not found in any audited source:
\begin{enumerate}
\item Richardson extrapolation of the \emph{logarithm} of truncated lacunary sinc products, followed by exponentiation to a multiplicative approximation;
\item the exact Gaussian-$q$-binomial formula for those log-extrapolation weights together with the full transformed Bernoulli--zeta tail;
\item the sign theorem and resulting alternating certified brackets for the infinite product;
\item the exact bounded condition number $(-q;q)_r/(q;q)_r$ of the level-$r$ log transform;
\item the Bell-polynomial relative-error and derivative hierarchy of the accelerated product;
\item transfer of this accelerated approximation to finite signed Thue--Morse exponential sums;
\item Lambert phase-locked Richardson nodes $x_j=(\lambda+j)2^{-(\lambda+j)}$ designed specifically to extract the periodic endpoint term $\Psi$;
\item exact transformed inverse-power moments on that lattice, including their complete-homogeneous and Bell-polynomial form.
\end{enumerate}

A targeted literature search located general Richardson-extrapolation theory and the standard ingredients separately, but no exact counterpart of the combined sinc-tail construction. The responsible statement is therefore ``new relative to the audited repository and apparently not standard,'' not a claim of global priority.

\subsection{Proof-status convention}

Every substantial assertion is tagged implicitly by its placement:
\begin{itemize}
\item Results in Sections~\ref{sec:tail}--\ref{sec:thue-morse-transfer} are \statusproved.
\item The phase-locked Lambert transform in Section~\ref{sec:lambert} is \statusconditional:
the interpolation algebra is proved here, while the all-orders endpoint input is inherited
from the formally verified endpoint theory.
\item Items in Section~\ref{sec:open} are explicitly conjectural or open.
\end{itemize}
Here ``proved'' means proved in this mathematical report, not necessarily formalized in
Lean.  In particular, the endpoint input is represented by
\texttt{fabiusSaddleLogCoefficient},
\path{log_fabius_sub_sharpLambertMain_hasAsymptoticExpansion}, and
\path{log_fabius_sub_sharpLambertExpansion_isBigO} in the Lean development,
whereas the sinc-tail and phase-extraction transforms remain formalization obligations.

\section{Notation and structural background}
\label{sec:notation}

\subsection{Sinc products and the Rvachev transform}

We use the entire normalized sinc function
\[
\sinc z=\begin{cases}\sin z/z,&z\ne0,\\1,&z=0.\end{cases}
\]
For $b>1$, define
\begin{equation}
P_b(t)=\prod_{k=1}^{\infty}\sinc\!\left(\frac{t}{b^k}\right),
\qquad
P_{b,N}(t)=\prod_{k=1}^{N}\sinc\!\left(\frac{t}{b^k}\right).
\label{eq:Pb-def}
\end{equation}
The dyadic case is $P=P_2$. In a common Fourier normalization for Rvachev's up-function,
\begin{equation}
\widehat{\Up}(\xi)=\prod_{h=0}^{\infty}\sinc\!\left(\frac{\pi\xi}{2^h}\right)
=\sinc(\pi\xi)P_2(\pi\xi).
\label{eq:up-factor}
\end{equation}
Thus every result for $P_2$ immediately applies to the nontrivial tail of $\widehat{\Up}$.

Set
\begin{equation}
q=b^{-2}\in(0,1),\qquad z_N=q^{N+1},\qquad
\rho_N=\frac{|t|}{b^{N+1}}=|t|z_N^{1/2}.
\label{eq:q-z-rho}
\end{equation}
The local analysis below requires $\rho_N<\pi$, exactly the condition that every omitted sinc factor lies before its first nonzero zero.

\subsection{\texorpdfstring{$q$}{q}-notation}

For $r\in\N$,
\[
(a;q)_r=\prod_{k=0}^{r-1}(1-aq^k),\qquad
(q;q)_r=\prod_{k=1}^{r}(1-q^k),
\]
and
\[
\qbinom{r}{j}{q}=\frac{(q;q)_r}{(q;q)_j(q;q)_{r-j}}
\quad(0\le j\le r).
\]
We will repeatedly use the finite $q$-binomial theorem
\begin{equation}
\prod_{k=0}^{r-1}(1+zq^k)
=\sum_{j=0}^{r}q^{\binom{j}{2}}\qbinom{r}{j}{q} z^j.
\label{eq:q-binomial}
\end{equation}

\subsection{Bernoulli and Bell polynomials}

The Bernoulli polynomials are defined by
\begin{equation}
\frac{te^{xt}}{e^t-1}=\sum_{n=0}^{\infty}B_n(x)\frac{t^n}{n!},
\qquad B_n=B_n(0).
\label{eq:bernoulli-gen}
\end{equation}
The complete exponential Bell polynomials $\mathcal B_n$ satisfy
\begin{equation}
\exp\!\left(\sum_{m\ge1}x_m\frac{u^m}{m!}\right)
=\sum_{n\ge0}\mathcal B_n(x_1,\ldots,x_n)\frac{u^n}{n!}.
\label{eq:bell-gen}
\end{equation}
They are the natural conversion device from logarithmic/cumulant data to multiplicative/moment data.

\subsection{The lower-Lambert coordinate}

Let $L=\log 2$. For sufficiently small $x>0$, set
\begin{equation}
\lambda(x)=-\frac{1}{L}W_{-1}(-Lx).
\label{eq:lambda-def}
\end{equation}
Then
\begin{equation}
x=\lambda 2^{-\lambda},\qquad \lambda>\frac1L.
\label{eq:lambda-inverse}
\end{equation}
The lower branch is essential: on $\lambda>1/L$, the map $\lambda\mapsto\lambda2^{-\lambda}$ is strictly decreasing and hence invertible.

\part{Exact logarithmic tails and multiplicative \texorpdfstring{$q$}{q}-Richardson acceleration}

\section{The exact Bernoulli--zeta tail of a lacunary sinc product}
\label{sec:tail}

\subsection{The logarithmic sinc series}

\begin{lemma}[Bernoulli--zeta expansion]
\label{lem:logsinc}
For $|u|<\pi$,
\begin{equation}
\log\frac{\sin u}{u}
=-\sum_{m=1}^{\infty}\gamma_m u^{2m},
\qquad
\gamma_m=\frac{\zeta(2m)}{m\pi^{2m}}
=\frac{2^{2m-1}|B_{2m}|}{m(2m)!}.
\label{eq:logsinc}
\end{equation}
Every $\gamma_m$ is strictly positive.
\end{lemma}

\begin{proof}
Euler's product
\(
\sin u/u=\prod_{n\ge1}(1-u^2/(\pi^2n^2))
\)
is locally uniformly convergent for $|u|<\pi$. Expanding $\log(1-v)=-\sum_{m\ge1}v^m/m$ and interchanging the absolutely convergent sums gives
\[
\log\frac{\sin u}{u}
=-\sum_{m\ge1}\frac{u^{2m}}{m\pi^{2m}}
\sum_{n\ge1}n^{-2m}.
\]
The Bernoulli form follows from Euler's evaluation of $\zeta(2m)$.
\end{proof}

The first coefficients are
\begin{equation}
\log\sinc u=-\frac{u^2}{6}-\frac{u^4}{180}
-\frac{u^6}{2835}-\frac{u^8}{37800}-\cdots.
\label{eq:logsinc-first}
\end{equation}

\subsection{Exact omitted-tail formula}

\begin{theorem}[Exact logarithmic tail]
\label{thm:exact-tail}
Let $b>1$, $N\ge0$, and $t\in\R$ satisfy $\rho_N<\pi$. If $P_{b,N}(t)\ne0$, then $P_b(t)$ and $P_{b,N}(t)$ have the same sign and
\begin{equation}
\log|P_{b,N}(t)|-\log|P_b(t)|
=\sum_{m=1}^{\infty}d_m(t)z_N^m,
\label{eq:tail-z}
\end{equation}
where
\begin{equation}
d_m(t)=\frac{\gamma_m t^{2m}}{1-q^m}
=\frac{\zeta(2m)t^{2m}}{m\pi^{2m}(1-q^m)}.
\label{eq:dm}
\end{equation}
Equivalently,
\begin{equation}
\log|P_{b,N}(t)|-\log|P_b(t)|
=\sum_{m=1}^{\infty}
\frac{\zeta(2m)}{m\pi^{2m}(1-q^m)}\rho_N^{2m}.
\label{eq:tail-rho}
\end{equation}
The series is absolutely convergent and strictly positive for $t\ne0$.
\end{theorem}

\begin{proof}
For $k\ge N+1$ we have $|t|/b^k\le\rho_N<\pi$, so every omitted sinc factor is positive.
Moreover, the logarithms of the omitted factors are absolutely summable by
\cref{lem:logsinc}; their product therefore converges to a strictly positive, nonzero
limit.  Hence truncation and limit have the same sign.  Applying the same expansion and
using absolute convergence,
\begin{align*}
\log|P_b(t)|-\log|P_{b,N}(t)|
&=\sum_{k=N+1}^{\infty}\log\sinc(t/b^k)\\
&=-\sum_{m=1}^{\infty}\gamma_m t^{2m}
\sum_{k=N+1}^{\infty}b^{-2mk}\\
&=-\sum_{m=1}^{\infty}\gamma_m t^{2m}
\frac{q^{m(N+1)}}{1-q^m}.
\end{align*}
Negating yields the result.
\end{proof}

\begin{corollary}[Analytic Bernoulli correction]
\label{cor:analytic-correction}
For $M\ge0$, define
\begin{equation}
C_{b,N}^{[M]}(t)=P_{b,N}(t)
\exp\!\left(-\sum_{m=1}^{M}d_m(t)z_N^m\right).
\label{eq:analytic-corrected}
\end{equation}
Then $C_{b,N}^{[M]}(t)$ has the same sign as $P_b(t)$ and
\begin{equation}
\log\frac{|C_{b,N}^{[M]}(t)|}{|P_b(t)|}
=\sum_{m=M+1}^{\infty}d_m(t)z_N^m>0
\quad(t\ne0).
\label{eq:analytic-correction-error}
\end{equation}
Thus $|C_{b,N}^{[M]}(t)|$ is a rigorous upper bound for $|P_b(t)|$.
\end{corollary}

This correction is explicit but uses $\zeta(2m)$ or Bernoulli numbers. The next construction achieves cancellation from product values alone.

\section{Geometric Lagrange weights and their \texorpdfstring{$q$}{q}-binomial form}
\label{sec:qweights}

\subsection{Interpolation at \texorpdfstring{$1,q,\ldots,q^r$}{1, q, ..., q-to-the-r}}

For an integer $r\ge0$, let $\lambda_{r,j}(q)$ be the Lagrange basis weight for evaluation at $0$ using nodes $q^0,\ldots,q^r$:
\begin{equation}
\lambda_{r,j}(q)=\prod_{\substack{0\le k\le r\\k\ne j}}
\frac{-q^k}{q^j-q^k}.
\label{eq:lambda-def-q}
\end{equation}

\begin{theorem}[Gaussian-$q$-binomial weights]
\label{thm:qweights}
For $0\le j\le r$,
\begin{equation}
\boxed{
\lambda_{r,j}(q)
=\frac{(-1)^{r-j}q^{\binom{r-j+1}{2}}}
{(q;q)_j(q;q)_{r-j}}
=\frac{(-1)^{r-j}q^{\binom{r-j+1}{2}}}{(q;q)_r}
\qbinom{r}{j}{q}.}
\label{eq:qweights}
\end{equation}
Their power moments are
\begin{equation}
S_{r,m}(q):=\sum_{j=0}^{r}\lambda_{r,j}(q)q^{jm}
=\frac{q^{rm}(q^{1-m};q)_r}{(q;q)_r}.
\label{eq:qmoment-full}
\end{equation}
Consequently,
\begin{equation}
S_{r,0}=1,
\qquad S_{r,m}=0\quad(1\le m\le r),
\label{eq:qmoment-cancel}
\end{equation}
and, for $m\ge r+1$,
\begin{equation}
S_{r,m}(q)
=(-1)^r q^{r(r+1)/2}
\frac{(q^{m-r};q)_r}{(q;q)_r}.
\label{eq:qmoment-positive}
\end{equation}
In particular, every nonzero residual moment has sign $(-1)^r$.
\end{theorem}

\begin{proof}
Splitting the product in \cref{eq:lambda-def-q} into $k<j$ and $k>j$ gives \cref{eq:qweights}. For the moments, substitute the Gaussian-binomial form and change variables $k=r-j$. The finite $q$-binomial theorem \cref{eq:q-binomial}, with a suitable monomial factor extracted, gives \cref{eq:qmoment-full}. The zeros for $1\le m\le r$ arise because $(q^{1-m};q)_r$ contains the factor $1-q^0$. For $m\ge r+1$, reverse the negative powers in the product to obtain \cref{eq:qmoment-positive}.
\end{proof}

\subsection{An exact condition number}

The alternating weights might appear dangerous, but their total variation is exceptionally well behaved.

\begin{theorem}[Bounded $\ell^1$ norm]
\label{thm:qcondition}
The absolute weight sum is
\begin{equation}
\kappa_r(q):=\sum_{j=0}^{r}|\lambda_{r,j}(q)|
=\frac{(-q;q)_r}{(q;q)_r}
=\prod_{k=1}^{r}\frac{1+q^k}{1-q^k}.
\label{eq:qcondition}
\end{equation}
Hence
\begin{equation}
\kappa_r(q)\nearrow \kappa_\infty(q)
:=\frac{(-q;q)_\infty}{(q;q)_\infty}<\infty.
\label{eq:qcondition-limit}
\end{equation}
For $q=1/4$,
\begin{equation}
\kappa_\infty(1/4)=1.9692603536682694319\ldots.
\label{eq:qcondition-dyadic}
\end{equation}
\end{theorem}

\begin{proof}
Set $k=r-j$ in the absolute-value sum and use \cref{eq:qweights}:
\[
\kappa_r(q)=\frac1{(q;q)_r}
\sum_{k=0}^{r}q^{k(k+1)/2}\qbinom{r}{k}{q}.
\]
\Cref{eq:q-binomial} with $z=q$ identifies the numerator as
$\prod_{i=0}^{r-1}(1+q^{i+1})=(-q;q)_r$. The infinite product converges because $\sum q^k<\infty$.
\end{proof}

\begin{remark}
This boundedness is not generic for high-order polynomial extrapolation. It is a special consequence of the geometric nodes and is one reason to extrapolate the logarithms: the inputs are of comparable scale, while the transformation itself has uniformly bounded absolute sensitivity.
\end{remark}

\section{Multiplicative \texorpdfstring{$q$}{q}-Richardson extrapolation}
\label{sec:qrichardson}

\subsection{Definition and exact transformed error}

Assume $\rho_N<\pi$ and $P_{b,N}(t)\ne0$. Define
\begin{equation}
\widehat P_{b;N,r}(t)
=\sgn(P_{b,N}(t))
\exp\!\left(\sum_{j=0}^{r}\lambda_{r,j}(q)
\log|P_{b,N+j}(t)|\right).
\label{eq:Phat}
\end{equation}
Equivalently,
\begin{equation}
|\widehat P_{b;N,r}(t)|
=\prod_{j=0}^{r}|P_{b,N+j}(t)|^{\lambda_{r,j}(q)}.
\label{eq:Phat-product}
\end{equation}
Negative exponents are harmless because the relevant truncations are nonzero.

\begin{theorem}[Exact multiplicative extrapolation error]
\label{thm:exact-rich-error}
Under the preceding hypotheses,
\begin{align}
E_{b;N,r}(t)
&:=\log\frac{|\widehat P_{b;N,r}(t)|}{|P_b(t)|}
\label{eq:E-def}\\
&=\sum_{m=r+1}^{\infty}d_m(t)z_N^mS_{r,m}(q)
\label{eq:E-series}\\
&=\frac{(-1)^rq^{r(r+1)/2}}{(q;q)_r}
\sum_{m=r+1}^{\infty}
\frac{\zeta(2m)}{m\pi^{2m}}
\frac{(q^{m-r};q)_r}{1-q^m}\rho_N^{2m}.
\label{eq:E-explicit}
\end{align}
The series is absolutely convergent. If $t\ne0$, then
\begin{equation}
\sgn E_{b;N,r}(t)=(-1)^r.
\label{eq:E-sign}
\end{equation}
\end{theorem}

\begin{proof}
Apply \cref{thm:exact-tail} to each truncation $N+j$:
\[
\log|P_{b,N+j}|=\log|P_b|+
\sum_{m\ge1}d_m z_N^m q^{jm}.
\]
The constant term survives because $S_{r,0}=1$, while the first $r$ powers vanish by \cref{eq:qmoment-cancel}. This proves \cref{eq:E-series}. \Cref{eq:E-explicit} follows from \cref{eq:qmoment-positive}. Every factor in its sum is positive, proving the sign statement.
\end{proof}

\subsection{Certified alternating brackets}

\begin{corollary}[Rigorous upper and lower bounds]
\label{cor:brackets}
For $t\ne0$ under the hypotheses of \cref{thm:exact-rich-error},
\begin{equation}
|\widehat P_{b;N,2s+1}(t)|<|P_b(t)|<
|\widehat P_{b;N,2s}(t)|
\qquad(s\ge0).
\label{eq:bracket}
\end{equation}
Thus any two consecutive extrapolation levels give a certified enclosure of the infinite product. In particular,
\begin{equation}
\left|\log\frac{|\widehat P_{b;N,r}(t)|}{|P_b(t)|}\right|
\le
\left|\log\frac{|\widehat P_{b;N,r}(t)|}
{|\widehat P_{b;N,r+1}(t)|}\right|
\label{eq:aposteriori-log}
\end{equation}
whenever the two levels lie on opposite sides as in \cref{eq:bracket}.
\end{corollary}

The bracket is stronger than a mere asymptotic claim: it follows from the exact sign of every term in the transformed tail.

\subsection{Leading term and explicit remainder bound}

\begin{corollary}[Sharp leading error]
\label{cor:leading-error}
As $\rho_N\to0$ with $r$ fixed,
\begin{equation}
E_{b;N,r}(t)
=C_r(q)\rho_N^{2r+2}+O(\rho_N^{2r+4}),
\label{eq:leading-error}
\end{equation}
where
\begin{equation}
C_r(q)=(-1)^r q^{r(r+1)/2}
\frac{\zeta(2r+2)}{(r+1)\pi^{2r+2}(1-q^{r+1})}.
\label{eq:Cr}
\end{equation}
For the dyadic case $q=1/4$,
\begin{equation}
C_0=\frac29,\qquad
C_1=-\frac1{675},\qquad
C_2=\frac1{178605},\qquad
C_3=-\frac1{154224000}.
\label{eq:Cr-dyadic}
\end{equation}
\end{corollary}

\begin{theorem}[Nonasymptotic error bound]
\label{thm:error-bound}
Let $a=(\rho_N/\pi)^2<1$. Then
\begin{equation}
|E_{b;N,r}(t)|
\le
\frac{q^{r(r+1)/2}\zeta(2)}
{(q;q)_r(1-q)(r+1)}
\frac{a^{r+1}}{1-a}.
\label{eq:error-bound}
\end{equation}
Consequently, for every fixed admissible pair $(t,N)$,
\begin{equation}
\widehat P_{b;N,r}(t)\longrightarrow P_b(t)
\qquad(r\to\infty)
\label{eq:r-convergence}
\end{equation}
with a supergeometric factor $q^{r(r+1)/2}$.
\end{theorem}

\begin{proof}
For $m\ge r+1$, \cref{eq:qmoment-positive} gives
\[
|S_{r,m}(q)|\le\frac{q^{r(r+1)/2}}{(q;q)_r}.
\]
Also $\zeta(2m)\le\zeta(2)$, $m^{-1}\le(r+1)^{-1}$, and $(1-q^m)^{-1}\le(1-q)^{-1}$. Sum the resulting geometric series in $a$. Since $(q;q)_r$ tends to the positive limit $(q;q)_\infty$, the right-hand side tends to zero faster than $q^{r^2/2}$ times an exponential factor.
\end{proof}

\subsection{Richardson table and multiplicative recurrence}

The transform can be evaluated without forming the closed weights.

\begin{algorithm}[Logarithmic $q$-Richardson table]
\label{alg:rich-table}
Set
\[
R_{N,0}(t)=\log|P_{b,N}(t)|.
\]
For $s\ge1$, recursively define
\begin{equation}
R_{N,s}(t)=
\frac{R_{N+1,s-1}(t)-q^sR_{N,s-1}(t)}{1-q^s}.
\label{eq:rich-recursion}
\end{equation}
Then
\begin{equation}
R_{N,r}(t)=\sum_{j=0}^{r}\lambda_{r,j}(q)
\log|P_{b,N+j}(t)|,
\qquad
\widehat P_{b;N,r}=\sgn(P_{b,N})e^{R_{N,r}}.
\label{eq:rich-table-result}
\end{equation}
\end{algorithm}

In multiplicative form,
\begin{equation}
|\widehat P_{b;N,s}|
=\left(
\frac{|\widehat P_{b;N+1,s-1}|}
{|\widehat P_{b;N,s-1}|^{q^s}}
\right)^{1/(1-q^s)}.
\label{eq:multiplicative-recursion}
\end{equation}
The logarithmic implementation is preferable because it avoids underflow in the rapidly decaying Fourier regime.

\subsection{Explicit dyadic levels}

For $b=2$, $q=1/4$. The first transforms are
\begin{align}
|\widehat P_{N,1}|
&=|P_N|^{-1/3}|P_{N+1}|^{4/3},
\label{eq:r1-dyadic}\\
|\widehat P_{N,2}|
&=|P_N|^{1/45}|P_{N+1}|^{-4/9}|P_{N+2}|^{64/45},
\label{eq:r2-dyadic}\\
|\widehat P_{N,3}|
&=|P_N|^{-1/2835}|P_{N+1}|^{4/135}
|P_{N+2}|^{-64/135}|P_{N+3}|^{4096/2835}.
\label{eq:r3-dyadic}
\end{align}
The log errors begin respectively at orders $\rho_N^4$, $\rho_N^6$, and $\rho_N^8$.

\subsection{Hybrid Bernoulli--Richardson correction}

The exact tail coefficients and the data-driven transform can be combined.

\begin{corollary}[Hybrid correction]
\label{cor:hybrid}
For $s\ge0$, define
\begin{equation}
\widehat P_{b;N,r}^{[s]}(t)
=\widehat P_{b;N,r}(t)
\exp\!\left(-\sum_{m=r+1}^{r+s}
 d_m(t)z_N^mS_{r,m}(q)\right).
\label{eq:hybrid}
\end{equation}
Then
\begin{equation}
\log\frac{|\widehat P_{b;N,r}^{[s]}(t)|}{|P_b(t)|}
=\sum_{m=r+s+1}^{\infty}d_m(t)z_N^mS_{r,m}(q),
\label{eq:hybrid-error}
\end{equation}
which still has sign $(-1)^r$. Thus the hybrid retains certified bracketing while moving the first nonzero term from degree $r+1$ to degree $r+s+1$.
\end{corollary}

This is the most direct point at which Bernoulli numbers and Gaussian $q$-binomial weights cooperate in one exact numerical scheme.

\section{Bell-polynomial error calculus and derivatives}
\label{sec:bell}

\subsection{From logarithmic error to relative error}

Write
\begin{equation}
e_m^{(r)}(t)=d_m(t)S_{r,m}(q),
\qquad e_m^{(r)}=0\quad(1\le m\le r).
\label{eq:em}
\end{equation}
Then \cref{eq:E-series} is $E=\sum_{m\ge1}e_m^{(r)}z_N^m$. Exponentiation and \cref{eq:bell-gen} give an exact convergent expansion.

\begin{theorem}[Bell-polynomial relative-error expansion]
\label{thm:bell-relative}
Under the hypotheses of \cref{thm:exact-rich-error},
\begin{equation}
\frac{|\widehat P_{b;N,r}(t)|}{|P_b(t)|}
=1+\sum_{n=r+1}^{\infty}
\frac{1}{n!}
\mathcal B_n\!\left(1!e_1^{(r)},2!e_2^{(r)},\ldots,n!e_n^{(r)}\right)z_N^n.
\label{eq:bell-relative}
\end{equation}
In particular, because $e_1^{(r)}=\cdots=e_r^{(r)}=0$, the first relative-error coefficient equals the first logarithmic-error coefficient:
\begin{equation}
\frac{|\widehat P_{b;N,r}|}{|P_b|}-1
=e_{r+1}^{(r)}z_N^{r+1}+O(z_N^{r+2}).
\label{eq:first-relative}
\end{equation}
\end{theorem}

The next coefficients exhibit the familiar cumulant-to-moment pattern. For example,
\begin{align*}
[z_N^{r+2}]\left(e^E-1\right)&=e_{r+2}^{(r)}
&& (r\ge1),\\
[z_N^{2r+2}]\left(e^E-1\right)&=e_{2r+2}^{(r)}+\frac12(e_{r+1}^{(r)})^2.
\end{align*}

\subsection{Logarithmic derivatives}

On an interval avoiding zeros, differentiation of \cref{eq:tail-z} is legitimate termwise. Let
\[
D^\ell=\frac{\dd^\ell}{\dd t^\ell},\qquad
(2m)_{\ell}=2m(2m-1)\cdots(2m-\ell+1).
\]

\begin{theorem}[Derivative acceleration]
\label{thm:derivative-accel}
For every fixed $\ell\ge0$ and admissible $t$ away from zeros,
\begin{equation}
D^\ell\!\left(
\sum_{j=0}^{r}\lambda_{r,j}\log|P_{b,N+j}(t)|
-\log|P_b(t)|\right)
=
\sum_{m\ge\max(r+1,\lceil\ell/2\rceil)}
\frac{\gamma_m(2m)_\ell t^{2m-\ell}}
{1-q^m}z_N^mS_{r,m}(q).
\label{eq:derivative-error}
\end{equation}
Thus the same $q$-Richardson table simultaneously cancels the first $r$ geometric tail powers in every logarithmic derivative.
\end{theorem}

For reconstructing derivatives of the product itself, set $L(t)=\log|P_b(t)|$. Fa\`a di Bruno's formula becomes
\begin{equation}
\frac{P_b^{(n)}(t)}{P_b(t)}
=\mathcal B_n\bigl(L'(t),L''(t),\ldots,L^{(n)}(t)\bigr)
\label{eq:bell-log-derivatives}
\end{equation}
where the sign of $P_b$ is locally constant. Combining \cref{eq:derivative-error,eq:bell-log-derivatives} gives an accelerated hierarchy for all derivatives of the Fourier transform away from its zeros.

\subsection{Application to interzero extrema}

Between consecutive zeros of $\widehat{\Up}$, an extremum of its modulus satisfies
\begin{equation}
\frac{\dd}{\dd\xi}\log|\widehat{\Up}(\xi)|=0.
\label{eq:peak-equation}
\end{equation}
\Cref{eq:up-factor} separates the first sinc factor from $P_2(\pi\xi)$. Theorem~\ref{thm:derivative-accel} gives a certified high-order approximation to the tail contribution in \cref{eq:peak-equation}. This suggests a numerically stable alternative to multiplying many tiny sinc factors and then differentiating the product.

\section{Transfer to finite signed Thue--Morse sums}
\label{sec:thue-morse-transfer}

Let
\begin{equation}
\varepsilon_n=(-1)^{s_2(n)},
\label{eq:epsilon}
\end{equation}
where $s_2(n)$ is the binary digit sum. The standard finite product identity
\begin{equation}
\sum_{n=0}^{2^N-1}\varepsilon_n z^n
=\prod_{j=0}^{N-1}(1-z^{2^j})
\label{eq:TM-poly}
\end{equation}
becomes a precise sinc identity after a phase-scaled substitution.

\begin{proposition}[Finite Thue--Morse representation of $P_N$]
\label{prop:TM-PN}
Define
\begin{equation}
T_N(t)=\sum_{n=0}^{2^N-1}\varepsilon_n
\exp\!\left(-\frac{\ii tn}{2^{N-1}}\right).
\label{eq:TN}
\end{equation}
Then
\begin{equation}
T_N(t)=\ii^N2^{-N(N-1)/2}
\e^{-\ii t(1-2^{-N})}t^N P_N(t),
\label{eq:TN-PN}
\end{equation}
with the identity at $t=0$ understood by continuity. Hence, for $t\ne0$,
\begin{equation}
P_N(t)=2^{N(N-1)/2}\ii^{-N}
\e^{\ii t(1-2^{-N})}t^{-N}T_N(t).
\label{eq:PN-TN}
\end{equation}
\end{proposition}

\begin{proof}
Insert $z=e^{-\ii t/2^{N-1}}$ in \cref{eq:TM-poly}. Reindex the factors to obtain
\[
T_N(t)=\prod_{k=0}^{N-1}(1-e^{-\ii t/2^k}).
\]
Using
\[
1-e^{-\ii a}=2\ii e^{-\ii a/2}\sin(a/2),
\]
sum the phases.  Finally, factor
\[
\prod_{j=1}^{N}\sin(t/2^j)
=t^N2^{-N(N+1)/2}P_N(t),
\]
which gives the result.
\end{proof}

\begin{corollary}[Accelerated infinite product from finite digital sums]
\label{cor:TM-accelerated}
For $t\ne0$ and admissible $N$,
\begin{equation}
|\widehat P_{N,r}(t)|
=\prod_{j=0}^{r}
\left(
2^{(N+j)(N+j-1)/2}|t|^{-(N+j)}|T_{N+j}(t)|
\right)^{\lambda_{r,j}(1/4)}.
\label{eq:TM-accelerated}
\end{equation}
The exact error, bracketing, and bounds of \cref{thm:exact-rich-error,cor:brackets,thm:error-bound} therefore become statements purely about a finite collection of signed Thue--Morse exponential sums.
\end{corollary}

This transfer is conceptually useful even when direct product evaluation is faster. It makes the approximation a finite digital identity and exposes a new route for symbolic or formal treatment through the already developed Thue--Morse polynomial machinery.

\section{Numerical verification and observed orders}
\label{sec:numerics}

All values in this section were recomputed with 100 decimal digits. The reference value used 300 sinc factors; at the displayed precision its omitted tail is negligible relative to every table entry. We report absolute logarithmic error
\[
\varepsilon_{N,r}(t)=
\left|\log|\widehat P_{N,r}(t)|-\log|P(t)|\right|,
\]
where $r=0$ denotes the raw truncation $P_N$.

\begin{table}[htbp]
\centering
\caption{Dyadic product at $t=3$: absolute logarithmic errors.}
\label{tab:numeric-t3}
\begin{tabular}{@{}r r r r r@{}}
\toprule
$N$ & raw $r=0$ & $r=1$ & $r=2$ & $r=3$ \\
\midrule
2 & $3.1368194\,10^{-2}$ & $2.9611726\,10^{-5}$ & $1.5785787\,10^{-8}$ & $2.5743474\,10^{-12}$ \\
3 & $7.8198398\,10^{-3}$ & $1.8359337\,10^{-6}$ & $2.4411880\,10^{-10}$ & $9.9424663\,10^{-15}$ \\
4 & $1.9535830\,10^{-3}$ & $1.1451700\,10^{-7}$ & $3.8045692\,10^{-12}$ & $3.8728157\,10^{-17}$ \\
5 & $4.8830986\,10^{-4}$ & $7.1537455\,10^{-9}$ & $5.9408270\,10^{-14}$ & $1.5117515\,10^{-19}$ \\
6 & $1.2207210\,10^{-4}$ & $4.4705340\,10^{-10}$ & $9.2810541\,10^{-16}$ & $5.9042379\,10^{-22}$ \\
\bottomrule
\end{tabular}
\end{table}

The successive-$N$ error ratios approach
\[
4^{-1},\quad4^{-2},\quad4^{-3},\quad4^{-4}
\]
for levels $r=0,1,2,3$, exactly as \cref{eq:leading-error} predicts. The signs also alternate: raw and level two are upper approximations in modulus; levels one and three are lower approximations.

\begin{table}[htbp]
\centering
\caption{Errors for two additional arguments, illustrating that the order is independent of $t$ inside the convergence disk.}
\label{tab:numeric-other}
\begin{tabular}{@{}c r r r r r@{}}
\toprule
$t$ & $N$ & raw $r=0$ & $r=1$ & $r=2$ & $r=3$ \\
\midrule
1 & 4 & $2.1701954\,10^{-4}$ & $1.4129551\,10^{-9}$ & $5.2149215\,10^{-15}$ & $5.8978483\,10^{-21}$ \\
1 & 6 & $1.3563390\,10^{-5}$ & $5.5189741\,10^{-12}$ & $1.2730608\,10^{-18}$ & $8.9985206\,10^{-26}$ \\
5 & 4 & $5.4288846\,10^{-3}$ & $8.8466438\,10^{-7}$ & $8.1669316\,10^{-11}$ & $2.3096377\,10^{-15}$ \\
5 & 6 & $3.3909800\,10^{-4}$ & $3.4497407\,10^{-9}$ & $1.9894410\,10^{-14}$ & $3.5155982\,10^{-20}$ \\
\bottomrule
\end{tabular}
\end{table}

For $t=3$ the product itself is
\begin{equation}
P(3)=0.585719176535068482912529262451\ldots.
\label{eq:P3}
\end{equation}
Even at $N=2$, the third transformed level improves the raw log error by approximately ten decimal orders while using only truncations through $N+3=5$.

\part{Lambert phase-locked extraction of the endpoint oscillation}

\section{The all-orders endpoint template}
\label{sec:lambert-template}

The repository's corrected endpoint analysis has the structural form
\begin{equation}
\log F(x)=\mathcal W_{\mathrm{MSE}}(x)+\Psi(\lambda(x))
+\sum_{m=1}^{M}\frac{A_m(\lambda(x))}{\lambda(x)^m}
+O\!\left(\lambda(x)^{-M-1}\right),
\label{eq:endpoint-template}
\end{equation}
where $\Psi$ and the coefficient functions $A_m$ are $1$-periodic. The nonoscillatory main saddle exponent may be written in the normalization used in the repository as
\begin{equation}
\mathcal W_{\mathrm{MSE}}(x)
=-\frac{L}{2}\lambda^2+
\left(1+\frac L2\right)\lambda
-\frac12\log\lambda+C_F,
\qquad L=\log2,
\label{eq:Wmse}
\end{equation}
with the convention-dependent constant $C_F$. The critical fact for the present construction is not the explicit constant but the periodicity of every $A_m$ in the same phase variable.

Define the reduced endpoint observable
\begin{equation}
Y(x)=\log F(x)-\mathcal W_{\mathrm{MSE}}(x).
\label{eq:Y}
\end{equation}
Then \cref{eq:endpoint-template} becomes
\begin{equation}
Y(x)=\Psi(\lambda)+\sum_{m\ge1}A_m(\lambda)\lambda^{-m}
\label{eq:Y-expansion}
\end{equation}
in the asymptotic sense.

\section{Phase-locked Lambert nodes}
\label{sec:lambert}

\subsection{Exact phase locking}

Fix $\lambda>1/L$ and define
\begin{equation}
\lambda_j=\lambda+j,
\qquad
x_j=x_j(\lambda)=(\lambda+j)2^{-(\lambda+j)},
\qquad j=0,1,\ldots.
\label{eq:phase-nodes}
\end{equation}

\begin{lemma}[Exact Lambert inversion on the lattice]
\label{lem:phase-lock}
For every $j\ge0$,
\begin{equation}
\lambda(x_j)=\lambda+j.
\label{eq:phase-inverse}
\end{equation}
Consequently, for every $1$-periodic function $G$,
\begin{equation}
G(\lambda(x_j))=G(\lambda).
\label{eq:periodic-lock}
\end{equation}
\end{lemma}

\begin{proof}
\Cref{eq:lambda-inverse} gives the inverse relation exactly on the decreasing branch. Periodicity gives the second assertion.
\end{proof}

Thus the nodes are not merely close in phase; they have exactly identical phase. Their physical spacing is nongeometric:
\begin{equation}
\frac{x_{j+1}}{x_j}
=\frac{\lambda+j+1}{2(\lambda+j)},
\qquad
\log_2\frac{x_j}{x_0}
=\log_2\frac{\lambda+j}{\lambda}-j.
\label{eq:x-spacing}
\end{equation}
This slight deviation from a dyadic ray is precisely what preserves the Lambert phase.

\subsection{Lagrange weights in reciprocal scale}

Set
\begin{equation}
h_j=\frac1{\lambda+j}.
\label{eq:h-nodes}
\end{equation}
We interpolate $Y(x_j)$ as a formal/asymptotic power series in $h_j$ and evaluate at $h=0$.

\begin{theorem}[Phase-locked Richardson weights]
\label{thm:lambert-weights}
The Lagrange weights for evaluation at $h=0$ from $h_0,\ldots,h_r$ are
\begin{equation}
\boxed{
\omega_{r,j}(\lambda)
=\prod_{\substack{0\le k\le r\\k\ne j}}
\frac{-h_k}{h_j-h_k}
=\frac{(-1)^{r-j}}{r!}\binom{r}{j}(\lambda+j)^r.}
\label{eq:omega}
\end{equation}
They satisfy
\begin{equation}
\sum_{j=0}^{r}\omega_{r,j}=1,
\qquad
\sum_{j=0}^{r}\omega_{r,j}(\lambda+j)^{-m}=0
\quad(1\le m\le r).
\label{eq:omega-cancel}
\end{equation}
\end{theorem}

\begin{proof}
Since
\[
h_j-h_k=\frac{k-j}{(\lambda+j)(\lambda+k)},
\]
each factor in the Lagrange product equals $(\lambda+j)/(j-k)$. Therefore
\[
\omega_{r,j}=\frac{(\lambda+j)^r}
{\prod_{k\ne j}(j-k)}
=\frac{(-1)^{r-j}}{j!(r-j)!}(\lambda+j)^r.
\]
The moment identities follow either from polynomial interpolation or from the fact that an $r$th finite difference annihilates every polynomial of degree below $r$.
\end{proof}

The first three levels are
\begin{align}
(\omega_{1,0},\omega_{1,1})
&=(-\lambda,\lambda+1),
\label{eq:omega1}\\
(\omega_{2,0},\omega_{2,1},\omega_{2,2})
&=\left(\frac{\lambda^2}{2},-(\lambda+1)^2,
\frac{(\lambda+2)^2}{2}\right),
\label{eq:omega2}\\
(\omega_{3,0},\omega_{3,1},\omega_{3,2},\omega_{3,3})
&=\left(-\frac{\lambda^3}{6},\frac{(\lambda+1)^3}{2},
-\frac{(\lambda+2)^3}{2},\frac{(\lambda+3)^3}{6}\right).
\label{eq:omega3}
\end{align}

\subsection{All residual moments}

Let $h_n(a_0,\ldots,a_r)$ denote the complete homogeneous symmetric polynomial, characterized by
\begin{equation}
\prod_{j=0}^{r}(1-a_jz)^{-1}
=\sum_{n=0}^{\infty}h_n(a_0,\ldots,a_r)z^n.
\label{eq:complete-homogeneous}
\end{equation}

\begin{theorem}[Exact inverse-power moments]
\label{thm:lambert-moments}
For $s\ge1$,
\begin{equation}
\sum_{j=0}^{r}\omega_{r,j}(\lambda+j)^{-(r+s)}
=\frac{(-1)^r}{(\lambda)_{r+1}}
 h_{s-1}\!\left(\frac1\lambda,\frac1{\lambda+1},\ldots,
\frac1{\lambda+r}\right),
\label{eq:lambert-moments}
\end{equation}
where $(\lambda)_{r+1}=\lambda(\lambda+1)\cdots(\lambda+r)$ is the rising factorial. In particular,
\begin{equation}
\sum_{j=0}^{r}\omega_{r,j}(\lambda+j)^{-(r+1)}
=\frac{(-1)^r}{(\lambda)_{r+1}}.
\label{eq:lambert-leading-moment}
\end{equation}
\end{theorem}

\begin{proof}
Substituting \cref{eq:omega} reduces the left-hand side to
\[
\frac1{r!}\sum_{j=0}^{r}(-1)^{r-j}\binom{r}{j}(\lambda+j)^{-s}.
\]
For $s=1$, partial fractions of $1/(\lambda)_{r+1}$ give \cref{eq:lambert-leading-moment}. Repeated differentiation with respect to $\lambda$, or multiplication of the generating function of reciprocal powers, yields \cref{eq:lambert-moments}. Equivalently, compare coefficients in
\[
\frac{1}{(\lambda-z)_{r+1}}
=\frac{1}{(\lambda)_{r+1}}
\prod_{j=0}^{r}\left(1-\frac{z}{\lambda+j}\right)^{-1}.
\]
\end{proof}

Define generalized harmonic sums
\begin{equation}
H_k^{(r)}(\lambda)=\sum_{j=0}^{r}(\lambda+j)^{-k}.
\label{eq:Hrk}
\end{equation}
Then the complete homogeneous polynomial in \cref{eq:lambert-moments} has the Bell representation
\begin{equation}
h_n\!\left(\frac1\lambda,\ldots,\frac1{\lambda+r}\right)
=\frac1{n!}\mathcal B_n\!\left(
0!H_1^{(r)},1!H_2^{(r)},\ldots,(n-1)!H_n^{(r)}
\right).
\label{eq:h-Bell}
\end{equation}
This is the second exact Bell-polynomial bridge in the report: the first exponentiated a log error; the second resolves all residual moments of the Lambert extrapolator.

\subsection{Extraction of the periodic endpoint factor}

Define
\begin{equation}
\widehat\Psi_r(\lambda)
=\sum_{j=0}^{r}\omega_{r,j}(\lambda)Y(x_j(\lambda)).
\label{eq:Psi-hat}
\end{equation}

\begin{theorem}[Phase-locked extraction of $\Psi$]
\label{thm:Psi-extraction}
Assume the all-orders expansion \cref{eq:endpoint-template} holds uniformly along a fixed compact phase set. Then, as $\lambda\to\infty$ with its fractional part fixed,
\begin{align}
\widehat\Psi_r(\lambda)-\Psi(\lambda)
&\sim \frac{(-1)^r}{(\lambda)_{r+1}}
\sum_{s=0}^{\infty}A_{r+1+s}(\lambda)
 h_s\!\left(\frac1\lambda,\ldots,\frac1{\lambda+r}\right)
\label{eq:Psi-full-error}\\
&=\frac{(-1)^rA_{r+1}(\lambda)}{(\lambda)_{r+1}}
+O(\lambda^{-r-2}).
\label{eq:Psi-leading-error}
\end{align}
In particular, the first $r$ inverse-$\lambda$ corrections are cancelled exactly.
\end{theorem}

\begin{proof}
At the phase-locked nodes, periodicity gives
\[
Y(x_j)=\Psi(\lambda)+\sum_{m\ge1}A_m(\lambda)(\lambda+j)^{-m}.
\]
The constant survives because the weights sum to one. Terms $1\le m\le r$ vanish by \cref{eq:omega-cancel}. Write $m=r+1+s$ and apply \cref{thm:lambert-moments}. A small uniformity point is essential: the absolute weight sum is $O_r(\lambda^r)$. To obtain the expansion through residual index $s=S$, apply the input all-orders theorem through order $2r+S+1$. Its transformed remainder is then $O_r(\lambda^{-r-S-2})$, while the finitely many undisplayed transformed coefficients start at the same order. This proves the displayed Poincar\'e expansion despite the growing weights.
\end{proof}

Exponentiating produces a multiplicative estimator of the periodic amplitude:
\begin{equation}
\widehat{\mathcal A}_r(\lambda)
=\exp(\widehat\Psi_r(\lambda))
=\prod_{j=0}^{r}
\left(\frac{F(x_j)}{\exp(\mathcal W_{\mathrm{MSE}}(x_j))}
\right)^{\omega_{r,j}(\lambda)}.
\label{eq:Ahat}
\end{equation}
Its relative-error coefficients are again complete Bell polynomials applied to the logarithmic coefficients in \cref{eq:Psi-full-error}.

\subsection{Coefficient extraction beyond \texorpdfstring{$\Psi$}{Psi}}

The same nodes can estimate $A_1,A_2,\ldots$ by differentiating the interpolating polynomial at $h=0$. If $\ell_j(h)$ is the Lagrange basis for $h_j$, then
\begin{equation}
\widehat A_s^{(r)}(\lambda)
=\frac1{s!}\sum_{j=0}^{r}\ell_j^{(s)}(0)Y(x_j),
\qquad 0\le s\le r,
\label{eq:As-extractor}
\end{equation}
with $\widehat A_0^{(r)}=\widehat\Psi_r$. For $s=1$,
\begin{equation}
\ell_j'(0)
=-\omega_{r,j}\sum_{\substack{0\le k\le r\\k\ne j}}(\lambda+k).
\label{eq:first-derivative-weight}
\end{equation}
Higher derivative weights are elementary symmetric polynomials in the numbers $-(\lambda+k)$ and can be generated by Bell polynomials. Thus a single phase-locked sample block can recover a triangular set of periodic coefficient functions.

\section{Bernoulli-polynomial geometry and conditioning of the Lambert lattice}
\label{sec:lambert-conditioning}

Unlike the geometric $q$-weights, the phase-locked weights grow polynomially in $\lambda$:
\begin{equation}
\kappa_r^{\mathrm L}(\lambda)
:=\sum_{j=0}^{r}|\omega_{r,j}(\lambda)|
=\frac1{r!}\sum_{j=0}^{r}\binom{r}{j}(\lambda+j)^r.
\label{eq:lambert-condition}
\end{equation}
Equivalently, if $J\sim\mathrm{Binomial}(r,1/2)$,
\begin{equation}
\kappa_r^{\mathrm L}(\lambda)
=\frac{2^r}{r!}\mathbb E(\lambda+J)^r
\sim_{\lambda\to\infty,\,r\,\mathrm{fixed}}\frac{2^r}{r!}\lambda^r.
\label{eq:lambert-condition-asymp}
\end{equation}
Thus high precision is essential. This is not a defect of the derivation: extrapolating from $h_j\asymp1/\lambda$ to the remote point $h=0$ is intrinsically ill-conditioned.

The binomial moments can be expanded through Stirling numbers:
\begin{equation}
\sum_{j=0}^{r}\binom{r}{j}j^k
=\sum_{\ell=0}^{k}
\left\{\begin{matrix}k\\\ell\end{matrix}\right\}
r^{\underline{\ell}}\,2^{r-\ell}.
\label{eq:binomial-moments}
\end{equation}
Here $r^{\underline{\ell}}=r(r-1)\cdots(r-\ell+1)$ is a falling factorial. For unweighted arithmetic-node diagnostics, Bernoulli polynomials give the exact power sums
\begin{equation}
\sum_{j=0}^{r}(\lambda+j)^p
=\frac{B_{p+1}(\lambda+r+1)-B_{p+1}(\lambda)}{p+1}.
\label{eq:bernoulli-power-sum}
\end{equation}
\Cref{eq:bernoulli-power-sum} is useful for scaling, centering, and constructing orthogonalized polynomial bases on the Lambert node block. It also makes explicit why Bernoulli polynomials are the natural arithmetic-lattice companion to the Bernoulli numbers already present in the sinc-tail coefficients.

\subsection{A stable evaluation strategy}

For practical extraction of $\Psi$:
\begin{enumerate}
\item compute $\lambda$ and $x_j$ at precision exceeding the target by roughly
$\log_{10}\kappa_r^{\mathrm L}(\lambda)$ guard digits;
\item evaluate $Y(x_j)$ with the same phase convention and main term;
\item use barycentric interpolation in $h_j$ rather than explicitly summing large alternating monomials;
\item repeat at successive $\lambda$ with the same fractional part and verify the predicted $\lambda^{-(r+1)}$ decay;
\item compare $r$ and $r+1$ transforms as an a posteriori diagnostic, while remembering that no sign bracket is presently available because the signs of $A_m(\lambda)$ are not known globally.
\end{enumerate}

\part{Unified interpretation, formalization path, and research program}

\section{Two extrapolation geometries}
\label{sec:dual-geometries}

The two constructions are instances of one principle: choose a coordinate in which the defect has a power expansion, then interpolate to the zero-defect point.

\begin{table}[htbp]
\centering
\caption{Comparison of the two extrapolation geometries.}
\label{tab:geometries}
\begin{tabularx}{\textwidth}{@{}l X X@{}}
\toprule
Feature & Sinc-tail transform & Lambert endpoint transform \\
\midrule
Defect coordinate & $z_N=q^{N+1}$ & $h=1/\lambda$ \\
Node family & $z_Nq^j$ & $(\lambda+j)^{-1}$ \\
Invariant structure & lacunary tail coefficients & periodic phase coefficients \\
Interpolation algebra & Gaussian $q$-binomial & ordinary binomial/finite difference \\
Leading cancellation & powers $z,\ldots,z^r$ & powers $h,\ldots,h^r$ \\
Coefficient source & Bernoulli numbers / $\zeta(2m)$ & periodic saddle jets $A_m$ \\
Reconstruction & exponential Bell polynomials & complete homogeneous and Bell polynomials \\
Conditioning & uniformly bounded in $r$ & grows like $2^r\lambda^r/r!$ \\
Certification & exact alternating brackets & asymptotic order, no global sign yet \\
\bottomrule
\end{tabularx}
\end{table}

The sinc scheme is unusually favorable because the geometric-node Lagrange weights have bounded total variation. The Lambert scheme pays for exact phase locking with growing weights, but it solves a different problem: extracting a nonconstant periodic amplitude that ordinary dyadic extrapolation would mix across phases.

\section{Consequences for Fabius and Rvachev computations}
\label{sec:consequences}

\subsection{Certified Fourier evaluation}

Using \cref{eq:up-factor}, compute the first factor exactly and bracket $P_2(\pi\xi)$ by consecutive transformed levels. This produces a certified enclosure for $|\widehat\Up(\xi)|$ away from zeros. Because the work is done in logarithms, the method remains usable deep in the log-normal decay regime where direct products underflow.

\subsection{Peak localization and envelope testing}

Use the derivative transform in \cref{eq:derivative-error} inside \cref{eq:peak-equation}. The resulting roots can test fine envelope conjectures without contaminating the comparison by truncation error. In particular, a claimed subleading Fourier envelope should be compared against a product value whose numerical error is below the proposed correction term; the bracket theorem makes that requirement checkable.

\subsection{Independent extraction of the endpoint oscillation}

The phase-locked transform \cref{eq:Psi-hat} gives a second way to recover $\Psi$: not from its Fourier/Gamma--zeta representation, but from $F$ itself. Agreement between the two routes would be a stringent validation of constants, branches, and phase conventions in the endpoint asymptotic.

\subsection{A bridge from exact dyadic values to asymptotics}

The repository already relates exact dyadic values to a phase-locked Lambert coordinate on fixed-numerator rays. The new construction suggests a complementary experiment:
\begin{enumerate}
\item use exact or high-precision values of $F$ at nearby phase-locked non-dyadic nodes $x_j$;
\item extract $\Psi$ and several $A_m$;
\item compare them with the periodic coefficient functions inferred from exact dyadic sequences whose phases become dense;
\item use the discrepancy to isolate the dyadic-arithmetic contribution from the universal saddle contribution.
\end{enumerate}
This does not by itself prove the exact-to-asymptotic bridge, but it supplies a sharply conditioned diagnostic coordinate for that problem.

\subsection{Base-\texorpdfstring{$b$}{b} generalized up-functions}

Theorems~\ref{thm:exact-tail}--\ref{thm:error-bound} hold for every real $b>1$. If a compactly supported refinement function has Fourier transform
\(
\prod_{k\ge1}\sinc(t/b^k)
\)
or a finite product times this tail, the same acceleration and bracketing apply with $q=b^{-2}$. This broadens the method beyond the dyadic Rvachev function and may be useful for nonbinary infinite convolutions.

\section{Lean formalization blueprint}
\label{sec:lean}

The sinc-side results are well suited to formalization because their analytic and algebraic components separate cleanly.

\subsection{Suggested definitions}

\begin{enumerate}
\item \texttt{lacunarySincProd b N t}: finite product over $k=1,\ldots,N$.
\item \texttt{lacunarySincProdInf b t}: infinite product, with convergence proved from $\sum t^2b^{-2k}$.
\item \texttt{qRichardsonWeight q r j}: the closed form \cref{eq:qweights}.
\item \texttt{qRichardsonLog q r f N}: weighted sum of $f(N+j)$.
\item \texttt{lambertPhaseNode lambda j}: $(\lambda+j)2^{-(\lambda+j)}$.
\item \texttt{lambertRichardsonWeight lambda r j}: \cref{eq:omega}.
\end{enumerate}

\subsection{Algebraic theorem chain}

The lowest-risk formalization order is:
\begin{enumerate}
\item prove the product form of $\lambda_{r,j}$ and its Gaussian-binomial simplification;
\item prove the moment identity \cref{eq:qmoment-full} from the finite $q$-binomial theorem;
\item prove the absolute weight sum \cref{eq:qcondition};
\item prove the Lambert weight formula and cancellation moments;
\item prove the reciprocal moment identity \cref{eq:lambert-moments} by rational-function coefficients;
\item identify the complete-homogeneous coefficients with the already
      formalized complete-Bell power-sum specialization.
\end{enumerate}
These steps require no delicate asymptotics.

\paragraph{Post-snapshot algebraic progress.}
The later module \texttt{GeometricLagrange.lean} now formalizes Lagrange
evaluation weights for arbitrary distinct nodes, polynomial and module-valued
exactness, the first omitted monomial moment, and geometric-node
specializations.  At the dyadic ratio \(q=1/2\), the later
\texttt{FabiusDiscreteLimitToeplitz.lean} identifies the closed
\(q\)-binomial row with those Lagrange weights and proves its all-degree
reverse moments and exact total variation.  These results do not yet supply
the arbitrary-ratio formulas used above---in particular the \(q=1/4\)
residual moments and condition number.  Independently,
\texttt{ExponentialPartition.lean} and
\texttt{ExponentialBell.lean} identify the universal weighted-partition
formula for an exponential coefficient with
\texttt{SaddleExpansion.expCoeff} over every commutative
\(\mathbb Q\)-algebra.  The later \texttt{MomentCumulantAlgebra.lean}
factorially normalizes the same engine into named complete-Bell and inverse
moment--cumulant transforms, with the classical Bell recurrence and
normalized inverse laws.  In particular, the coefficient mechanism behind
\eqref{eq:h-Bell} can reuse these bridges with
\(E_k=k^{-1}\sum_i a_i^k\) for \(k\ge1\).  The displayed
complete-homogeneous specialization, the Richardson error formulas, and every
analytic estimate in this report remain distinct formalization obligations.

\subsection{Analytic theorem chain}

The analytic side can then proceed as follows:
\begin{enumerate}
\item establish the log-sinc power series on $(-\pi,\pi)$;
\item justify the double-series interchange for the omitted tail;
\item derive the exact transformed error by finite linearity;
\item prove positivity of every factor in \cref{eq:E-explicit};
\item deduce bracketing and the explicit majorant;
\item add termwise differentiation on compact zero-free intervals.
\end{enumerate}

The endpoint input is already formalized with its exact phase and every finite remainder:
the periodic coefficient functions are
\texttt{fabiusSaddleLogCoefficient}, and the expansion is proved in
\path{FabiusFullAsymptoticExpansion.lean}.  What remains in the frontier namespace is the
phase-locked extraction transform itself and the verification that its weighted use of
those uniform remainders has precisely the strength asserted in
\cref{thm:Psi-extraction}.  The arbitrary-node interpolation foundation is already
available in \path{GeometricLagrange.lean}; the closed reciprocal-grid weights and their
all-residual moments are the next algebraic obligations.

\section{Open problems and next frontier steps}
\label{sec:open}

\subsection{Optimal order under finite precision}

For exact input logs, \cref{eq:error-bound} improves supergeometrically with $r$, while \cref{eq:qcondition} remains bounded. In floating-point arithmetic, however, the truncations $P_{N+j}$ are themselves computed recursively and share correlated rounding errors. A sharp forward-error theorem should combine:
\[
\text{analytic error }\asymp q^{r(r+1)/2}a^{r+1}
\quad\text{with}\quad
\text{roundoff }\lesssim\kappa_r(q)u.
\]
Because $\kappa_r(q)$ converges, the optimal $r$ may be much larger than in ordinary Richardson extrapolation.

\subsection{Monotonicity within each bracket parity}

The sign theorem places every even level above the limit in modulus and every odd level
below it, but it does not order two levels of the same parity.  In fact, such an ordering is
false over the full parameter range $b>1$.  To see this directly, take
\[
 q=\frac45,\qquad b=q^{-1/2}=\frac{\sqrt5}{2},\qquad N=0,
 \qquad \rho_0=\pi\sqrt{\frac{99}{100}},
 \qquad t=b\rho_0=\pi\sqrt{\frac{99}{80}}.
\]
This is admissible because $\rho_0<\pi$, and $P_{b,0}=1$ is the nonzero empty product.
Substitution in the absolutely convergent series \eqref{eq:E-explicit}, with
$a=(\rho_0/\pi)^2=99/100$, gives
\[
\begin{array}{c|rrrr}
r&0&1&2&3\\ \hline
E_{b;0,r}(t)
& 13.5220607386&-12.9320597704&16.5469385988&-14.7912158217.
\end{array}
\]
Thus $E_{b;0,2}>E_{b;0,0}>0$ and
$E_{b;0,3}<E_{b;0,1}<0$: both proposed same-parity monotonicities fail.
These numbers are reproducible by summing the terms through $m=5000$ in
\eqref{eq:E-explicit}.  The same majorization used in \cref{thm:error-bound} bounds the
discarded tail, for each $0\le r\le3$, by
\[
 \frac{q^{r(r+1)/2}\zeta(2)}
 {(q;q)_r(1-q)(5001)}\frac{a^{5001}}{1-a}<2\times10^{-22}.
\]

The correctly parameterized open question is therefore: for which pairs
$(q,a)\in(0,1)^2$, and for which levels $s$, do
\[
 E_{2s}(q,a)>E_{2s+2}(q,a)>0,
 \qquad
 E_{2s+1}(q,a)<E_{2s+3}(q,a)<0,
\]
where $E_r(q,a)$ denotes the right-hand side of \eqref{eq:E-explicit} after replacing
$(\rho_N/\pi)^2$ by $a$?  No nonempty parameter regime for these additional inequalities
is asserted here.  Whenever they do hold, consecutive same-parity levels sharpen the
certified brackets of \cref{cor:brackets}.

\subsection{Global continuation across zeros}

The logarithmic construction is naturally local to a zero-free component. A global formulation could factor all explicitly known zeros from the truncations and extrapolate the analytic nonvanishing remainder. Such a theorem would give uniform certified evaluation across large frequency intervals, including neighborhoods of high-multiplicity dyadic zeros.

\subsection{Direct formulas for the phase coefficients}

The transformed moments \cref{eq:lambert-moments} show exactly how $A_{r+1},A_{r+2},\ldots$ enter the extraction error. This invites an inverse problem: compute $\widehat\Psi_r$ and $\widehat A_s^{(r)}$ at many phases, then identify their Fourier coefficients in terms of the Gamma--zeta data already present in the endpoint theory.

\subsection{A sign theory for Lambert corrections}

A sign pattern for $A_m(\lambda)$ would turn the phase-locked transform into a certified bracket analogous to \cref{cor:brackets}. Even a sign theorem for $A_1$ or $A_2$ on a substantial phase interval would be useful.

\subsection{Thue--Morse arithmetic of transformed products}

For $q=1/4$, every exponent in \cref{eq:r1-dyadic,eq:r2-dyadic,eq:r3-dyadic} is rational. Raising to a common denominator converts the approximant to a rational monomial identity in finite Thue--Morse sums. It is natural to ask whether the numerators and denominators exhibit systematic $2$-adic cancellation analogous to the exact dyadic arithmetic of $F$.

\section{Conclusion}

The repository already contains the ingredients of a remarkable analytic-combinatorial machine: Thue--Morse cancellation, Gaussian interpolation, exact dyadic arithmetic, infinite sinc products, Bernoulli cumulants, Bell reconstruction, and Lambert-$W$ endpoint geometry. The new results here arise by changing \emph{what} is extrapolated and \emph{where} the nodes are placed.

For the Fourier product, extrapolating logarithms exposes an exact positive power series in the geometric tail coordinate. Gaussian-$q$-binomial weights annihilate its first powers, and the remaining $q$-moments all have one sign. This yields high-order multiplicative approximants, exact alternating brackets, explicit bounds, bounded conditioning, Bell-polynomial relative errors, derivative acceleration, and a finite Thue--Morse realization.

For the endpoint asymptotic, moving by integer steps in the exact lower-Lambert coordinate freezes every periodic coefficient. Lagrange interpolation in reciprocal $\lambda$ then isolates the periodic factor and its correction hierarchy. The residual moments are not opaque: they are complete homogeneous symmetric polynomials, hence Bell polynomials in generalized harmonic sums, while Bernoulli polynomials govern the arithmetic-node geometry.

The two mechanisms are complementary. One is globally algebraic in the truncation index and exceptionally stable; the other is phase-sensitive and designed to separate a periodic asymptotic invariant. Both are sufficiently explicit to support immediate symbolic checking, high-precision computation, and a staged Lean formalization.

\appendix

\section{Original audited TeX inventory}
\label{app:inventory}

This appendix records the 57 distinct \texttt{.tex} paths included in the original
repository-wide overlap audit.  It is provenance for that audit, not a current corpus
index: subsequent frontier reports increased the present tree to 67 TeX sources.
Historical files are frozen versions and were treated as prior art even when their content
has since been consolidated elsewhere.

\subsection{Living top-level and canonical sources (6)}
\begin{enumerate}
\item \path{docs/FabiusFunction_Mathematical_Glossary/FabiusFunction_Mathematical_Glossary.tex}
\item \path{docs/Fabius_Function_and_Rvachev_Up/Fabius_Function_and_Rvachev_Up.tex}
\item \path{docs/Non_Elementarity_of_the_Fabius_Function/Non_Elementarity_of_the_Fabius_Function.tex}
\item \path{docs/fabius_lean_walkthrough/fabius_lean_walkthrough.tex}
\item \path{docs/non-formalized-research-frontiers/non-formalized-research-frontiers.tex}
\item \path{docs/non-formalized-research-frontiers/drafts/Thue_Morse_Formula_Atlas/Thue_Morse_Formula_Atlas.tex}
\end{enumerate}

\subsection{Living Rvachev Fourier-decay drafts (12)}
\begin{enumerate}[resume]
\item \path{docs/non-formalized-research-frontiers/drafts/rvachev_up_fourier_decay/Rvachev_Up_Fourier_Decay-2/Rvachev_Up_Fourier_Decay.tex}
\item \path{docs/non-formalized-research-frontiers/drafts/rvachev_up_fourier_decay/Rvachev_Up_Fourier_Decay-3/Fourier_decay_of_Rvachev_up.tex}
\item \path{docs/non-formalized-research-frontiers/drafts/rvachev_up_fourier_decay/Rvachev_Up_Fourier_Decay_Comparative_Audit/Rvachev_Up_Fourier_Decay_Comparative_Audit.tex}
\item \path{docs/non-formalized-research-frontiers/drafts/rvachev_up_fourier_decay/Rvachev_Up_Fourier_Decay_Gentle_Guide/Rvachev_Up_Fourier_Decay_Gentle_Guide.tex}
\item \path{docs/non-formalized-research-frontiers/drafts/rvachev_up_fourier_decay/Rvachev_Up_Fourier_Decay_Second_Wave_Audit/Rvachev_Up_Fourier_Decay_Second_Wave_Audit.tex}
\item \path{docs/non-formalized-research-frontiers/drafts/rvachev_up_fourier_decay/asymptotic-decay-rate-of-an-infinite-product-of-sinc-functions/asymptotic-decay-rate-of-an-infinite-product-of-sinc-functions.tex}
\item \path{docs/non-formalized-research-frontiers/drafts/rvachev_up_fourier_decay/rvachev_up_fourier_decay-1/rvachev_up_fourier_decay.tex}
\item \path{docs/non-formalized-research-frontiers/drafts/rvachev_up_fourier_decay/rvachev_up_fourier_decay-4/rvachev_up_fourier_decay.tex}
\item \path{docs/non-formalized-research-frontiers/drafts/rvachev_up_fourier_decay/rvachev_up_fourier_decay-5/rvachev_up_fourier_decay.tex}
\item \path{docs/non-formalized-research-frontiers/drafts/rvachev_up_fourier_decay/rvachev_up_fourier_decay-6/rvachev_up_fourier_decay.tex}
\item \path{docs/non-formalized-research-frontiers/drafts/rvachev_up_fourier_decay/rvachev_up_fourier_decay-7/rvachev_up_fourier_decay.tex}
\item \path{docs/non-formalized-research-frontiers/drafts/rvachev_up_fourier_decay/rvachev_up_fourier_decay-8/rvachev_up_fourier_decay.tex}
\end{enumerate}

\subsection{Bundled source-paper TeX (7)}
\begin{enumerate}[resume]
\item \path{docs/papers/AIF_2017__67_2_637_0/AIF_2017__67_2_637_0-corrected.tex}
\item \path{docs/papers/AIF_2017__67_2_637_0/AIF_2017__67_2_637_0.tex}
\item \path{docs/papers/arXiv-1502.06738v2/Lacunary_products.tex}
\item \path{docs/papers/arXiv-1702.05442v1/09-Function.tex}
\item \path{docs/papers/arXiv-1702.06487v3/157-Arithmetic-v3.tex}
\item \path{docs/papers/arXiv-2211.13570v1/document.tex}
\item \path{docs/papers/ubiq15/ubiq15-corrected.tex}
\end{enumerate}

\subsection{Frozen archive: early notes and standalone studies (5)}
\begin{enumerate}[resume]
\item \path{docs/archive/early-notes/Fabius Asymptotic/Fabius Asymptotic.tex}
\item \path{docs/archive/early-notes/K-fold summation over the signed Thue-Morse sequence/K-fold summation over the signed Thue-Morse sequence.tex}
\item \path{docs/archive/standalone-studies/Non_Elementarity_of_the_Fabius_Function/Non_Elementarity_of_the_Fabius_Function.tex}
\item \path{docs/archive/standalone-studies/Small_Argument_Asymptotics/Small_Argument_Asymptotics.tex}
\item \path{docs/archive/standalone-studies/fabius-inverse-dyadic-closed-form/fabius-inverse-dyadic-closed-form.tex}
\end{enumerate}

\subsection{Frozen archive: primary expositions and supplements (8)}
\begin{enumerate}[resume]
\item \path{docs/archive/primary-exposition/Fabius_Function_and_Rvachev_Up/Fabius_Function_and_Rvachev_Up.tex}
\item \path{docs/archive/primary-exposition/Fabius_Function_and_Rvachev_Up-2/Fabius_Function_and_Rvachev_Up.tex}
\item \path{docs/archive/primary-exposition/Fabius_Function_and_Rvachev_Up-3/Fabius_Function_and_Rvachev_Up.tex}
\item \path{docs/archive/primary-exposition/Fabius_and_Rvachev_Up-2/Fabius_and_Rvachev_Up.tex}
\item \path{docs/archive/primary-exposition/Fabius_and_Rvachev_up/Fabius_and_Rvachev_up.tex}
\item \path{docs/archive/primary-exposition/Fabius_function_and_up_function/Fabius_function_and_up_function.tex}
\item \path{docs/archive/primary-supplements/Fabius_Function_Missing_Parts/Fabius_Function_Missing_Parts.tex}
\item \path{docs/archive/primary-supplements/Fabius_Function_Missing_Parts-2/Fabius_Function_Missing_Parts.tex}
\end{enumerate}

\subsection{Frozen archive: q-calculus, walkthrough, and glossary (6)}
\begin{enumerate}[resume]
\item \path{docs/archive/q-calculus/Fabius_q_Special_Functions/Fabius_q_Special_Functions.tex}
\item \path{docs/archive/q-calculus/fabius-q-special-functions/fabius-q-special-functions.tex}
\item \path{docs/archive/lean-walkthrough/Fabius_Function_Lean_Walkthrough/Fabius_Function_Lean_Walkthrough.tex}
\item \path{docs/archive/lean-walkthrough/fabius_lean_walkthrough/fabius_lean_walkthrough.tex}
\item \path{docs/archive/glossary/FabiusFunction_Mathematical_Glossary/FabiusFunction_Mathematical_Glossary.tex}
\item \path{docs/archive/glossary/FabiusFunction_Mathematical_Glossary-2/FabiusFunction_Mathematical_Glossary.tex}
\end{enumerate}

\subsection{Frozen archive: research frontiers (13)}
\begin{enumerate}[resume]
\item \path{docs/archive/research-frontiers/Fabius_Dyadic_Asymptotic_Bridge/Fabius_Dyadic_Asymptotic_Bridge.tex}
\item \path{docs/archive/research-frontiers/Fabius_Dyadic_Formulae_and_Alternative_Representations/Fabius_Dyadic_Formulae_and_Alternative_Representations.tex}
\item \path{docs/archive/research-frontiers/Fabius_Dyadic_Formulae_to_Asymptotics/Fabius_Dyadic_Formulae_to_Asymptotics.tex}
\item \path{docs/archive/research-frontiers/Fabius_Dyadic_q_Connections/Fabius_Dyadic_q_Connections.tex}
\item \path{docs/archive/research-frontiers/Fabius_Thue_Morse_Convergence_Rate/Fabius_Thue_Morse_Convergence_Rate.tex}
\item \path{docs/archive/research-frontiers/Fabius_Thue_Morse_Convergence_Rates/Fabius_Thue_Morse_Convergence_Rates.tex}
\item \path{docs/archive/research-frontiers/Primary_Exposition_Gap_Register/Primary_Exposition_Gap_Register.tex}
\item \path{docs/archive/research-frontiers/Repeated_Integration_and_Rvachev_Up/Repeated_Integration_and_Rvachev_Up.tex}
\item \path{docs/archive/research-frontiers/Rvachev_Up_from_Repeated_Integration/Rvachev_Up_from_Repeated_Integration.tex}
\item \path{docs/archive/research-frontiers/non-formalized-research-frontiers/non-formalized-research-frontiers.tex}
\item \path{docs/archive/research-frontiers/Thue_Morse_Formula_Atlas-1/Thue_Morse_Formula_Atlas.tex}
\item \path{docs/archive/research-frontiers/Thue_Morse_Formula_Atlas-2/Consolidated_Formula_Atlas_for_the_Thue_Morse_Sequence.tex}
\item \path{docs/archive/research-frontiers/Thue_Morse_Formula_Atlas-3/Thue_Morse_Formula_Atlas.tex}
\end{enumerate}

\section{Additional algebraic details}
\label{app:details}

\subsection{Derivation of the \texorpdfstring{$q$}{q}-moment formula}

Starting from \cref{eq:qweights},
\begin{align*}
S_{r,m}(q)
&=\frac{1}{(q;q)_r}
\sum_{j=0}^{r}(-1)^{r-j}
q^{\binom{r-j+1}{2}}\qbinom{r}{j}{q} q^{jm}\\
&=\frac{q^{rm}}{(q;q)_r}
\sum_{k=0}^{r}(-1)^k q^{\binom{k+1}{2}-km}\qbinom{r}{k}{q}\\
&=\frac{q^{rm}}{(q;q)_r}
\sum_{k=0}^{r}q^{\binom{k}{2}}\qbinom{r}{k}{q}
(-q^{1-m})^k\\
&=\frac{q^{rm}(q^{1-m};q)_r}{(q;q)_r}.
\end{align*}
For $m\ge r+1$,
\begin{align*}
(q^{1-m};q)_r
&=\prod_{k=0}^{r-1}(1-q^{1-m+k})\\
&=(-1)^r q^{\sum_{k=0}^{r-1}(1-m+k)}
\prod_{k=0}^{r-1}(1-q^{m-1-k})\\
&=(-1)^r q^{r(1-m)+r(r-1)/2}(q^{m-r};q)_r.
\end{align*}
Multiplication by $q^{rm}$ gives the exponent $r(r+1)/2$ in \cref{eq:qmoment-positive}.

\subsection{Why all transformed sinc-tail terms have one sign}

The sign theorem is not a generic property of Richardson extrapolation. It relies on two positivity facts:
\begin{enumerate}
\item the log-sinc tail coefficients $d_m(t)$ are positive because $\zeta(2m)>0$;
\item after cancellation, every geometric moment is
\[
S_{r,m}(q)=(-1)^r\times
\frac{q^{r(r+1)/2}(q^{m-r};q)_r}{(q;q)_r},
\]
whose remaining factors are positive for $m\ge r+1$.
\end{enumerate}
Their product therefore has the same sign for every residual degree. This is the mechanism behind the certified brackets.

\subsection{Bell form of complete homogeneous functions}

Taking logarithms of \cref{eq:complete-homogeneous},
\[
\log\prod_{j=0}^{r}(1-a_jz)^{-1}
=\sum_{k=1}^{\infty}\frac{z^k}{k}\sum_{j=0}^{r}a_j^k.
\]
Comparing this with \cref{eq:bell-gen} gives
\[
h_n(a_0,\ldots,a_r)
=\frac1{n!}\mathcal B_n\left(
0!p_1,1!p_2,\ldots,(n-1)!p_n\right),
\]
where $p_k=\sum_j a_j^k$. Substituting $a_j=(\lambda+j)^{-1}$ proves \cref{eq:h-Bell}.

\section{Reproducibility notes}
\label{app:repro}

The numerical checks used arbitrary-precision arithmetic with the following direct definitions:
\begin{align*}
P_N(t)&=\prod_{k=1}^{N}\frac{\sin(t/2^k)}{t/2^k},\\
P(t)&\approx\prod_{k=1}^{300}\frac{\sin(t/2^k)}{t/2^k},\\
\log|\widehat P_{N,r}(t)|
&=\sum_{j=0}^{r}\lambda_{r,j}(1/4)\log|P_{N+j}(t)|.
\end{align*}
At 100-digit working precision, the 300-factor reference tail is far below the displayed errors. Symbolic checks verified the moment identities through $r=8$ and multiple values of $m$ before the general proofs were written.

\begin{thebibliography}{99}

\bibitem{ProveItPrimary}
V. Reshetnikov,
\emph{The Fabius Function and Rvachev's Up-Function},
living primary exposition in the \repo{} repository,
\texttt{Analysis/FabiusFunction/docs/Fabius\_Function\_and\_Rvachev\_Up/} (accessed 27 August 2026).

\bibitem{ProveItFrontier}
V. Reshetnikov,
\emph{Non-formalized Fabius Research Frontiers},
canonical quarantined frontier volume in the \repo{} repository
(accessed 27 August 2026).

\bibitem{ProveItAtlas}
V. Reshetnikov,
\emph{Thue--Morse Formula Atlas},
living draft in the \repo{} repository
(accessed 27 August 2026).

\bibitem{ProveItDecay}
V. Reshetnikov et al.,
\emph{Rvachev Up Fourier Decay} dossier series,
twelve living TeX reports in the \repo{} repository
(accessed 27 August 2026).

\bibitem{Arias1982}
J. Arias de Reyna,
\emph{An infinitely differentiable function with compact support: definition and properties},
Rev. Real Acad. Cienc. Exactas F\'is. Nat. Madrid \textbf{76} (1982), 21--38;
arXiv:1702.05442.

\bibitem{Arias2018}
J. Arias de Reyna,
\emph{Arithmetic of the Fabius function},
Integers \textbf{18} (2018), Article A51;
arXiv:1702.06487.

\bibitem{Corless1996}
R. M. Corless, G. H. Gonnet, D. E. G. Hare, D. J. Jeffrey, and D. E. Knuth,
\emph{On the Lambert $W$ function},
Advances in Computational Mathematics \textbf{5} (1996), 329--359.

\bibitem{Sidi1988}
A. Sidi,
\emph{Generalizations of Richardson extrapolation with applications to numerical integration},
1988, doi:10.1007/978-3-0348-6398-8\_22.

\bibitem{Comtet1974}
L. Comtet,
\emph{Advanced Combinatorics},
D. Reidel, Dordrecht, 1974.

\bibitem{AlloucheShallit2003}
J.-P. Allouche and J. Shallit,
\emph{Automatic Sequences: Theory, Applications, Generalizations},
Cambridge University Press, 2003.

\bibitem{AlloucheLacunary}
J.-P. Allouche, M. Mend\`es France, and J. Peyri\`ere,
works on lacunary trigonometric products and the Thue--Morse sequence;
source TeX included in the repository as \texttt{Lacunary\_products.tex}.

\end{thebibliography}

\end{document}
